% (c) Danny Lee Eldridge · AEGIS Community License · DRAFT — not published
\documentclass[aps,prd,twocolumn,superscriptaddress,nofootinbib,showkeys]{revtex4-2}
\usepackage{amsmath,amssymb,hyperref}
\begin{document}
\title{Flavour at the End of Inflation: Koide's Form, Trimaximal Mixing and a Quantised $\bar\theta$ from the Two Zero-Energy Vacua of an $SL(2,\mathbb Z)$ Plateau}
\author{Danny Eldridge}
\email{dannyeldridge54@gmail.com}
\affiliation{Independent researcher, Oklahoma City, Oklahoma, USA}
\date{19 November 2026 (draft)}
\begin{abstract}
In a companion paper~\cite{P1}, inflation is driven by the running of the species scale along the type IIB axio-dilaton $\tau$, with potential $V_0\propto|\nabla\ln E(\tau,\tfrac32)|^2$. That potential vanishes at exactly two points of the fundamental domain, the $SL(2,\mathbb Z)$ fixed points $\rho=e^{i\pi/3}$ and $i$. We ask what the flavour sector looks like if its modular fields come to rest at these two vacua. We report the following results.
(i) Two-field simulations of the end of inflation show that the inflaton itself ends at $\rho$ for every transverse kick tested. The line it falls along has a transverse instability band ($0.96\lesssim y\lesssim1.9$), which makes a ``Y-fork'' towards $i$, but the inflaton crosses it too quickly to be diverted.
(ii) The residual $\mathbb Z_3$ at $\rho$ forces a circulant $\sqrt{M_\ell}$. Its eigenvalues are exactly the Koide--Brannen form $\sqrt{m_k}=a+2|b|\cos(\delta+2\pi k/3)$, and $Q=2/3$ is equivalent to equal norm in the singlet and doublet components. The data give amplitude $\sqrt2$ and $\delta=2/9=Q/3$ to within $1\sigma$. Minimal $A_4$ weight-2 constructions placed exactly at $\rho$ do not produce these numbers.
(iii) With charged leptons at $\rho$ and neutrinos at $i$, $A_4$ fixes the second PMNS column to $(1/3,1/3,1/3)$ (TM2). This frees $\theta_{13}$, which exact-$\rho$ tribimaximal mixing excludes at $38\sigma$, and gives $\sin^2\theta_{12}=0.3408$ ($+2.8\sigma$). For $\sin^2\theta_{23}=0.561$ or $0.470$ it predicts $\delta_{CP}=236^\circ$ or $125^\circ$, or $287^\circ$ or $74^\circ$, respectively.
(iv) At $\rho$ the $A_4$ quark determinant is the real number $27/8$. With CP-symmetric couplings $\bar\theta\in\{0,\pi\}$, but the CKM phase vanishes ($|J|<5\times10^{-15}$). The observed phase therefore needs a quark modulus at a generic point with anomaly-free weights, as in Feruglio, Strumia and Titov~\cite{FST}.
All computations were pre-registered. We state plainly which steps are principles and which are derived.
\end{abstract}
\keywords{flavour; modular symmetry; $A_4$; Koide formula; trimaximal mixing; strong CP; inflation; $SL(2,\mathbb Z)$ fixed points; pre-registration}
\maketitle

\section{Two vacua from one potential}
The species-gradient potential $V_0=V_\infty(\tfrac23)^2y^2|\nabla\ln E(\tau,\tfrac32)|^2$ vanishes wherever $\nabla E=0$. In the fundamental domain this happens at $\rho$ and $i$. Numerically $V_0(\rho)=1.6\times10^{-22}$ and $V_0(i)=2.2\times10^{-22}$, and both points are strict minima, with $V_0\approx8.3\times10^{-8}$ (at $\rho$) and $(2.1$--$42)\times10^{-8}$ (at $i$) at distance $10^{-3}$. The inflaton falls along $x=\tfrac12$, which is a valley near $\rho$ ($y<0.954$) and a ridge for $0.96\lesssim y\lesssim1.9$, where the low side is the line $x=0$ that leads to $i$. We call this the Y-fork. A two-field integration with the hyperbolic kinetic metric shows the inflaton crossing the band and settling at $\rho$ for transverse kicks $\delta x=0$--$3\times10^{-2}$. These kicks are far larger than its quantum spread. Inflation ends at $y\approx2.4$, so the CMB observables of Ref.~\cite{P1} are unchanged. Whether a slow modulus drains to $i$ or to $\rho$ is computed in Sec.~V.

\section{Charged leptons at $\rho$: Koide's form}
At $\rho$ the residual symmetry is $\mathbb Z_3=\langle ST\rangle$. A square-root mass matrix covariant under the cyclic $\mathbb Z_3$ is circulant, $\sqrt M=a\mathbb 1+bC+\bar bC^2$, with eigenvalues
\begin{equation}
\sqrt{m_k}=a+2|b|\cos(\delta+2\pi k/3),\qquad Q=\tfrac13+\tfrac23|b|^2/a^2 .
\end{equation}
Koide's relation $Q=2/3$ is therefore $|b|^2=a^2/2$: the singlet and doublet components of $\sqrt M$ carry equal norm. The algebra is known~\cite{Foot}. What is new here is that inflation selects the $\mathbb Z_3$ point dynamically. Using Belle~II $m_\tau=1777.09\pm0.14$~MeV~\cite{Belle} we find $Q=0.6666735\pm0.0000077$ ($+0.9\sigma$), amplitude $1.414228$, and $\delta=0.2222136\pm0.0000095$ against $2/9$ ($-0.9\sigma$). With $m_\tau$ as the only input, $A=\sqrt2$ and $\delta=2/9$ give $m_\mu=105.666$~MeV and $m_e=0.511029$~MeV. The data also satisfy $\delta=Q/3$ to within $1\sigma$. We do not derive $\sqrt2$ or $2/9$. A locked test of four minimal $A_4$ weight-2 constructions exactly at $\rho$ gives $Q=0.360$, $0.341$, $1/3$ and $1/3$; all fail. The fixed point fixes the \emph{form}, not the numbers.

\section{Neutrinos at $i$: TM2}
At $i$ the residual symmetry is $\mathbb Z_2=\langle S\rangle$. In the $A_4$ basis of Ref.~\cite{F17}, the overlap of the $S=+1$ eigenvector with the $ST$ eigenbasis is $(1/3,1/3,1/3)$. Only the $\nu_2$ assignment survives, so this is TM2. With $\sin^2\theta_{13}=0.02195$ (NuFIT~6.0~\cite{NuFIT}), it gives $\sin^2\theta_{12}=1/(3\cos^2\theta_{13})=0.3408$ against $0.307^{+0.012}_{-0.011}$, i.e.\ $+2.8\sigma$. The allowed range is $\sin^2\theta_{23}\in[0.395,0.605]$. Table~\ref{tab:d} lists the $\delta_{CP}$ predictions, which were written before any comparison with $\delta_{CP}$ data. If everything sat at $\rho$, the result would be tribimaximal mixing, which is excluded at $38\sigma$ by $\theta_{13}$.
\begin{table}[b]\caption{TM2 predictions ($\theta_{13}$ from data).}\label{tab:d}
\begin{ruledtabular}\begin{tabular}{lcc}
$\sin^2\theta_{23}$ & 0.561 & 0.470\\ $\delta_{CP}$ & $235.6^\circ$ or $124.9^\circ$ & $286.6^\circ$ or $73.9^\circ$\\
\end{tabular}\end{ruledtabular}\end{table}

\section{Quarks: quantised $\bar\theta$ and the CKM phase}
For the standard $A_4$ structure $\mathrm{diag}(\alpha,\beta,\gamma)\cdot[[Y_1,Y_3,Y_2],[Y_2,Y_1,Y_3],[Y_3,Y_2,Y_1]]$, the determinant at $\rho$ is $Y_1^3+Y_2^3+Y_3^3-3Y_1Y_2Y_3=27/8$, which is real. Both $\rho$ and $i$ are CP-symmetric points. With real couplings $\bar\theta\in\{0,\pi\}$, and $\bar\theta=0$ is a discrete sign choice. In the same setting, however, $|J|<5\times10^{-15}$ over $6000$ samples, against the observed $3.1\times10^{-5}$. The CKM phase therefore requires the mechanism of Ref.~\cite{FST}: anomaly-free quark weights $\sum(2k_Q+k_u+k_d)=0$, so that $\det M_q$ is a real constant, together with a quark modulus at a generic point, for example $\tau_q\approx\tfrac18+i$. In this construction the quark modulus is not one of the two vacua.

\section{Do slow fields reach $i$?}
Whether a slow modulus, which does not carry the inflaton's momentum, drains to $i$ or to $\rho$ under the gradient flow of $V_0$ is being computed under a pre-registered lock (PREREG\_BASINS). The criterion, fixed in advance, is that a fraction of hyperbolic area $f_i\ge0.25$ draining to $i$ makes a modulus tied off at $i$ dynamically natural. The result will be reported in a revised version. Until then the assignment of the neutrino modulus to $i$ is a postulate.

\section{What is assumed}
We assume the following. The species-gradient potential itself~\cite{P1}. Equipartition, $|b|^2=a^2/2$, together with $\delta=Q/3$. That the lepton moduli come to rest at the fixed points, with the neutrino modulus at $i$. The anomaly-free quark weights and the quark modulus at a generic point. None of these is derived here. The vacuum energy beyond zero and the origin of $\Lambda_{\rm obs}$ are not addressed.

\emph{Integrity note.} All computations were pre-registered (locks D1595--D1607). Fixes to code were logged before any rerun, and one invalid run is kept on record. The $\delta_{CP}$ values are predictions, to be scored against DUNE and Hyper-Kamiokande.

\begin{thebibliography}{99}
\bibitem{P1} D.~Eldridge, ``Inflation from the running of the species scale'' (draft, 2026).
\bibitem{FST} F.~Feruglio, A.~Strumia, A.~Titov, JHEP 07 (2023) 027, arXiv:2305.08908.
\bibitem{Foot} R.~Foot, hep-ph/9402242; Y.~Koide, Lett.~Nuovo Cim.~34 (1982) 201; C.~A.~Brannen (2006).
\bibitem{Belle} Belle II, arXiv:2305.19116.
\bibitem{F17} F.~Feruglio, arXiv:1706.08749.
\bibitem{NuFIT} I.~Esteban et al., NuFIT 6.0, arXiv:2410.05380.
\end{thebibliography}
\end{document}
